% year: תשפ"ו
% semester: א
% moed: ב
% number: 5ב
% points: 10
% categories: improper-integrals
% official_solution: yes
% source: exams/מבחנים/תשפו/sol B.pdf

\begin{question}
(10 נק') בדקו התכנסות או התבדרות של האינטגרל המוכלל הבא:
\[ \int_1^\infty \frac{\cos x}{x^2 - \ln x} \]
הערה: אתם יכולים להניח שהמכנה חיובי בתחום האינטגרציה.
\end{question}

\begin{hint}
האינטגרנד אינו בעל סימן קבוע. בדקו התכנסות בהחלט: חסמו $|f(x)| \le \frac{1}{x^2 - \ln x}$ והשוו (במבחן ההשוואה הגבולי) ל-$\frac{1}{x^2}$.
\end{hint}

\begin{solution}
נסמן $f(x) = \frac{\cos x}{x^2 - \ln x}$. המכנה חיובי ב-$[1, \infty)$ (כפי שמותר להניח; למעשה $\ln x \le x - 1 < x \le x^2$ עבור $x \ge 1$), ולכן $f$ רציפה שם והאינטגרל מוכלל רק בגלל הגבול האינסופי. $f$ מחליפה סימן אינסוף פעמים, לכן נבדוק \textbf{התכנסות בהחלט}.

\textbf{חסימה:} לכל $x \ge 1$,
\[ 0 \le |f(x)| = \frac{|\cos x|}{x^2 - \ln x} \le \frac{1}{x^2 - \ln x} =: h(x). \]

\textbf{התכנסות $\int_1^\infty h$:} נשווה ל-$g(x) = \frac{1}{x^2}$. שתי הפונקציות חיוביות, ו-
\[ \lim_{x\to\infty} \frac{g(x)}{h(x)} = \lim_{x\to\infty} \frac{x^2 - \ln x}{x^2} = \lim_{x\to\infty} \left( 1 - \frac{\ln x}{x^2} \right) = 1, \]
כי $\frac{\ln x}{x^2} \to 0$ (לופיטל: $\frac{1/x}{2x} \to 0$). הגבול סופי וחיובי, ולכן לפי מבחן ההשוואה הגבולי $\int_1^\infty h$ ו-$\int_1^\infty g$ מתכנסים או מתבדרים יחד. מכיוון ש-$\int_1^\infty \frac{dx}{x^\alpha}$ מתכנס עבור $\alpha > 1$, ובפרט $\int_1^\infty \frac{dx}{x^2} = 1$, גם $\int_1^\infty h$ מתכנס.

\textbf{מסקנה:} לפי מבחן ההשוואה (כי $0 \le |f| \le h$), $\int_1^\infty |f|$ מתכנס. אינטגרל מוכלל שמתכנס בהחלט מתכנס, ולכן
\[ \int_1^\infty \frac{\cos x}{x^2 - \ln x}\,dx \ \text{מתכנס (ואף בהחלט).} \]
\end{solution}
